\chapter{Formal Proofs}
\begin{quotation}\textit{
One of the biggest and deepest challenges in computing is that of
getting things right. Errors or oversights in code might lead to
crashed, incorrect behaviour or security loopholes. If some small
component in a large project has problems these may bubble up and
impact the whole. A fine example of that arose in 1994 where some
Intel Pentium processors returned incorrect results for the quotient
of certain very particular floating point numbers. Of course few people
would have performed calculations that happened to use the numeric
values involved, and yet fewer would have noticed anything amiss. But
despite the issue being of rather low practical relevance it dented
confidence in every calculation performed in that model of CPU. \\
Even extensive and careful testing of code tends just to be able to
reveal bugs but not to certify their absence, and there are many cases
where flaws that lead to software misbehavior or insecurity have remained
present in major widely-used code for decades.\\
One response to this is to consider having formal proofs of code
correctness. Setting these up first calls for totally precise explanations
of what would consitute acceptable behaviour, and even when this is
available it is typically quite challenging.
%Operational vs. Denotational semantics.
%Abstract interpretation.
%How are you sure your proof checker is not bug-ridden?
}\end{quotation}


\section{Correct rounding}
The discussion here is not going to explain how the results it cites have
been achieved. But it notes a case where a team has deployed program proof
techniques to deliver a basic capability used by millions in such a way
that it is as good and reliable as possible.

Pretty well all computer languages provide functions for computing the
elementary functions: $\sin$, $\cos$, $\exp$, $\log$ and a bunch more.
historically in most cases the implementers accepted that the results
they delivered were really close to the ideal ones -- only the lowest
bits of the floating point values could deviate from perfection. After
all an error on one part on $10^17$ is really unlikely to hurt the users.
Neverthless there is an ideal against which perfection can be judged.
Any floating point number presented to one of these functions really
stands for an integer multiplied or divided by some power of 2. For instance
on a standard computer the number you present as $(1/3)$ is liable
to be stored as $(1501199875790165/2^{52})$. So one can take a view that
when ``$(1/3)$'' is presented to some function that ideal answer to be
returned will be the floating point value closest to the value
of the function when used on $(1501199875790165/2^{52})$. It should be
obvious that achieving this risks being seriously complicated and expensive!
Furthermore if someone has developed code that they claim achieves this,
checking it is not going to be straightforward since there are $2^64$
distinct bit-patterns that represent floating point values\footnote{Some
stand for infinities or for ``not numbers'' rather than being cases that need
careful checking}. This is way too large a number to make exhaustive checking
feasible, and it is not at all obvious that if just a few inputs lead to
very slightly incorrect results that there is a sane way to identify them.

The CORE-MATH project (see \verb+https://core-math.gitlabpages.inria.fr/+
uses a system called GAPPA to mechanise proofs associated with their
numeric code, concerning themselves at a really detailed level with every
rounding operation and every place where errors could creep in. Seting a
goal of obtaining results where for every input their versions of library
functions will return precisely correctly rounded results, they show that
this level of accuracy can be achieved while having code that runs as
fast (and sometimes faster) than existing elementary functions libraries.
For some of their code they publish human readable explanations of all the
cases that their proofs of correctness need to consider, but behind those
written explanations there will lie significant computational support.
Some widely-distributed systems are incoporating this work so that all
users -- most of whom will not even think about least-significant bit
accuracy -- can benefit from it. The message from this is that formal
proof can apply to code that is truly widely used, but that using it will
often be a lot of very specialist work!

 